\section*{Tóm tắt}
\addcontentsline{toc}{section}{Tóm tắt}

Nội dung của báo cáo được tổ chức thành 6 chương chính, trình bày chi tiết quy trình xây dựng hệ thống từ giai đoạn khảo sát, phân tích đến thiết kế và định hướng phát triển, cụ thể như sau:

\textbf{Chương 1: Giới thiệu} \\
Chương này trình bày tổng quan về bối cảnh thực hiện đề tài tại Trường Đại học Bách khoa (ĐHQG-HCM), nêu rõ các vấn đề bất cập trong công tác quản lý hoạt động ngoại khóa hiện tại. Từ đó, xác định mục tiêu, phạm vi nghiên cứu và đề xuất hướng giải quyết thông qua việc xây dựng hệ thống phần mềm quản lý tập trung.

\textbf{Chương 2: Cơ sở lý thuyết và công nghệ} \\
Tập trung giới thiệu các nền tảng lý thuyết và công nghệ chủ chốt được áp dụng trong đồ án. Nội dung bao gồm các quy trình nghiệp vụ Đoàn - Hội và bộ công nghệ (Tech Stack) được lựa chọn để phát triển ứng dụng Web cũng như công nghệ AI tiên tiến (RAG) phục vụ cho tính năng Chatbot.

\textbf{Chương 3: Yêu cầu về hệ thống} \\
Phân tích chi tiết các yêu cầu đối với hệ thống phần mềm, bao gồm yêu cầu chức năng (các nghiệp vụ quản lý thành viên, hoạt động, tài chính, biểu mẫu...) và yêu cầu phi chức năng (bảo mật, hiệu năng, trải nghiệm người dùng), làm cơ sở cho các bước thiết kế tiếp theo.

\textbf{Chương 4: Mô hình hóa hệ thống} \\
Chương này mô tả hệ thống dưới góc độ kỹ thuật thông qua ngôn ngữ mô hình hóa UML. Các biểu đồ Use Case, Activity và Sequence được xây dựng nhằm làm rõ các tác nhân, quy trình nghiệp vụ và luồng tương tác dữ liệu bên trong hệ thống.

\textbf{Chương 5: Phân tích và Thiết kế hệ thống} \\
Đi sâu vào thiết kế kiến trúc tổng thể và cấu trúc dữ liệu (ERD) của hệ thống. Đặc biệt, chương này trình bày chi tiết về kiến trúc và cơ chế hoạt động của Chatbot AI sử dụng kỹ thuật RAG (Retrieval-Augmented Generation), từ khâu xử lý dữ liệu đầu vào (Vector Embedding) đến quy trình truy vấn và phản hồi thông tin tự động.

\textbf{Chương 6: Tổng kết} \\
Đánh giá lại những kết quả đạt được của đồ án so với mục tiêu ban đầu, nhìn nhận các ưu điểm và hạn chế còn tồn tại. Đồng thời, chương này cũng đề xuất lộ trình phát triển, nâng cấp và bảo trì hệ thống trong tương lai.

\clearpage